% Appendices (protocol, benchmarks, implementation)
\appendix

\section{Experiment Details}
\label{app:experiment-details}

\subsection{Benchmarks}
\label{app:benchmarks}
\label{app:dataset-rationale}
\label{app:dataset-descriptions}

\paragraph{Benchmark selection rationale.}
The selection of benchmark datasets is governed by three requirements of \textbf{Adaptive Agent
Routing (AAR)}:
\textbf{(1)}~coverage of a wide spectrum of query complexity, from lightweight factual lookups to
deeply compositional multi-step reasoning;
\textbf{(2)}~structural diversity that stress-tests the \textbf{Query Complexity Estimator (QCE)}
and the learned routing policy across distinct task types; and
\textbf{(3)}~evaluation of four agent strategies in
$\mathcal{A}=\{\texttt{raw},\texttt{cot},\texttt{react},\texttt{multiagent}\}$ under conditions
that are meaningful and representative of deployed agentic systems.

\paragraph{Coverage of query complexity.}
A central thesis of AAR is that not all queries warrant the same agent configuration.
To validate this claim, the benchmark suite must span the complexity range that QCE is designed to
discriminate.
Our \emph{executed} evaluation uses fixed-size subsamples per family
(Table~\ref{tab:deployment-sizes}) rather than full upstream corpora.
The \emph{routing benchmark} draws queries from \textbf{HotpotQA}, \textbf{MuSiQue}, and
\textbf{MATH} for training and regret-based evaluation (\S\ref{sec:experiments}).
\emph{Benchmark transfer} additionally runs all four strategies on \textbf{GAIA} and
\textbf{MMLU-Pro}, which are held out from routing-benchmark construction; HotpotQA, MuSiQue, and
MATH are re-evaluated with independent transfer queries ($N{=}965$ pooled).
Pooled utility uses $\lambda{=}25$ (Table~\ref{tab:pooled-utility}).

\begin{table}[h]
  \centering
  \footnotesize
  \caption{Executed evaluation subsamples per benchmark family (transfer evaluation).}
  \label{tab:deployment-sizes}
  \begin{tabular}{@{}lr@{}}
    \toprule
    Workload & Queries \\
    \midrule
    GAIA & 165 \\
    HotpotQA & 200 \\
    MuSiQue & 200 \\
    MMLU-Pro & 200 \\
    MATH & 200 \\
    \textbf{Pooled (micro)} & \textbf{965} \\
    \bottomrule
  \end{tabular}
\end{table}

\paragraph{MMLU-Pro~\citep{wang2024mmlupro}.}
MMLU-Pro is an enhanced Massive Multitask Language Understanding benchmark introduced as frontier
models approached 86--87\% on the original MMLU, reducing its usefulness for tracking progress and
exhibiting sensitivity to prompt phrasing and scoring.
MMLU-Pro integrates more challenging, reasoning-focused questions, expands answer choices from
four to ten (lowering random-guess success), and comprises over 12{,}000 curated questions across
14 domains (e.g., biology, business, chemistry, computer science, economics, engineering, health,
history, law, mathematics, philosophy, physics, and psychology).
Reported results show accuracy drops of roughly 16--33 percentage points relative to MMLU and
lower prompt sensitivity (on the order of 2\% vs.\ 4--5\% variation on MMLU).
Chain-of-thought models tend to outperform direct answering on MMLU-Pro, indicating that the
benchmark rewards reasoning rather than surface recall.
In our study, MMLU-Pro is a held-out transfer family for knowledge-intensive multiple-choice
routing.

\paragraph{MATH~\citep{hendrycks2021math}.}
MATH measures competition-style mathematical problem solving with 12{,}500 problems from contests
including AMC~10, AMC~12, and AIME.
Each problem is tagged with difficulty level 1--5 and one of seven subjects (prealgebra, algebra,
number theory, counting and probability, geometry, intermediate algebra, and precalculus), with
full step-by-step solutions provided.
Unlike benchmarks focused on short numeric answers, MATH requires complete multi-step derivations
across a wide difficulty range.
MATH is used in both routing-benchmark construction and transfer; in our results it often favors
chain-of-thought over tool-heavy strategies (Table~\ref{tab:live-benchmark}).

\paragraph{GAIA~\citep{mialon2023gaia}.}
GAIA evaluates general AI assistants on real-world tasks requiring reasoning, multimodal inputs,
web browsing, and tool use, with short verifiable answers.
The public release contains 466 questions at three difficulty levels; many require multiple steps,
tool use (e.g., browsers or code interpreters), and heterogeneous inputs (images, video,
spreadsheets).
Reported human accuracy is near 92\% versus roughly 15\% for early GPT-4 with plugins, illustrating
a large human--machine gap on everyday-style agentic tasks.
We evaluate a 165-query subsample under full strategy execution as a held-out, tool-intensive
transfer workload.

\paragraph{HotpotQA~\citep{yang2018hotpotqa}.}
HotpotQA is a large-scale Wikipedia multi-hop QA benchmark with 113K question--answer pairs that
require locating and reasoning over multiple supporting documents.
Questions are diverse and not constrained to a fixed knowledge schema; sentence-level supporting
facts enable supervised reasoning and explanation.
In the distractor setting, each instance includes two gold paragraphs and eight TF-IDF distractors
(ten paragraphs total); the full-wiki setting requires retrieval over Wikipedia.
HotpotQA stresses retrieval, evidence integration, and multi-step inference, and is used in both
the routing benchmark and transfer evaluation.

\paragraph{MuSiQue~\citep{trivedi2022musique}.}
MuSiQue targets a shortcut-exploitation weakness in prior multi-hop datasets by composing
\emph{connected} single-hop questions into 2--4 hop chains.
The MuSiQue-Ans collection contains 25K such questions and is reported to be substantially harder,
with a wider human--machine gap and roughly a 30-point F1 drop for single-hop models relative to
multi-hop systems.
MuSiQue-Full adds unanswerable contrast questions for stricter evaluation.
By structurally requiring genuine compositional reasoning, MuSiQue complements HotpotQA in both
routing-benchmark construction and transfer.

\begin{table}[h]
  \centering
  \footnotesize
  \caption{Pooled deployment utility (965 queries).}
  \label{tab:pooled-utility}
  \begin{tabular}{@{}lr@{}}
    \toprule
    Policy & Mean utility \\
    \midrule
    Proposed complexity--semantic & 0.440 \\
    Heuristic baseline & 0.458 \\
    \midrule
    Bootstrap 95\% CI (heur.\ $-$ proposed) & $[-0.0006,\,+0.035]$ \\
    Two-sided $p$ (approx.) & 0.059 \\
    \bottomrule
  \end{tabular}
\end{table}

\subsection{Routing Benchmark Construction}
\label{app:benchmark-construction}

\paragraph{Routing benchmark.}
The routing benchmark is a supervised routing dataset constructed from oracle executions of
four reasoning strategies.
Each instance contains query features $\phi(q)$, per-strategy exact-match outcomes
$r_a(q)$, execution costs $c_a(q)$, and derived routing targets (hard labels and cost-soft
$p^\star(\cdot \mid q)$).
It supports all regret-based evaluations in the main paper; the deployed router never
observes per-query oracle outcomes at inference.

The benchmark aggregates queries from HotpotQA, MuSiQue, and MATH into a stratified split:
train $808$, validation $100$, test $100$.
For each query, all four strategies are executed once offline; outcomes and costs are
recorded in a centralized oracle log.
Hard oracle labels and cost-soft targets $p^\star(\cdot \mid q)$ are derived from these
recorded outcomes (Eqs.~\ref{eq:hard-oracle} and~\ref{eq:cost-soft}).
Splits are stratified by source dataset; the test split is held out for all primary
routing metrics reported in \S\ref{sec:experiments}.
Soft labels $p^\star(\cdot \mid q)$ are precomputed per split with cost temperature
$\tau$ (Appendix~\ref{app:hyperparameters}).

\section{Feature Definitions}
\label{app:features}

Table~\ref{tab:feature-configs-full} summarizes the routing representations evaluated
in \S\ref{sec:experiments}, including a complexity-only ablation ($\mathbf{z}$ without $\psi$).
Repository identifiers are provided only for reproducibility with the released codebase.

\begin{table}[h]
  \centering
  \footnotesize
  \caption{Routing representations.}
  \label{tab:feature-configs-full}
  \begin{tabular}{@{}p{0.24\linewidth}p{0.38\linewidth}rrl@{}}
    \toprule
    Representation & Description & Dim.\ $m$ & Planner? & Repository ID \\
    \midrule
    Embedding-only &
      Semantic query embedding $\psi(q)$ only &
      16 & No & \texttt{emb\_only} \\
    Heuristic baseline &
      Lexical heuristics plus $\psi(q)$; no planner decomposition &
      50 & No & \texttt{heur\_emb} \\
    Complexity-only &
      Planner complexity $\mathbf{z}(q)$ only (ablation) &
      5 & Yes & \texttt{cvec5} \\
    \textbf{Complexity--semantic} &
      Planner complexity $\mathbf{z}(q)$ plus $\psi(q)$ (proposed) &
      21 & Yes & \texttt{cvec5\_emb} \\
    \bottomrule
  \end{tabular}
\end{table}

\paragraph{Complexity--semantic (proposed).}
Planner-derived $\mathbf{z}(q)\in\mathbb{R}^5$ concatenated with PCA-projected
$\psi(q)\in\mathbb{R}^{16}$.
QCE aggregates procedure-graph statistics into five complexity dimensions; TrainNorm caps
are fit on training plans only (Appendix~\ref{app:hyperparameters}).

\paragraph{Heuristic baseline.}
Lexical heuristic features plus $\psi(q)$; no LLM decomposition.

\paragraph{Embedding-only.}
$\psi(q)$ only.

\subsection{QCE Dimension Ablation (Leave-One-Out)}
\label{app:qce-dim-loo}

Table~\ref{tab:qce-dim-loo} reports leave-one-out ablations on the routing benchmark test split
($n{=}100$): four complexity dimensions plus $\psi(q)$, with one dimension withheld.
All rows use cost-soft HGBM; total regret includes QCE planner cost.

\begin{table}[h]
  \centering
  \footnotesize
  \caption{QCE leave-one-out ablation (four dims + $\psi$, held-out test).}
  \label{tab:qce-dim-loo}
  \begin{tabular}{@{}lrr@{}}
    \toprule
    Withheld dimension & Total regret $\downarrow$ & Agent regret $\downarrow$ \\
    \midrule
    \textbf{Full (none withheld)} & \textbf{0.198} & \textbf{0.192} \\
    Reasoning ($z_2$) & 0.199 & 0.193 \\
    Evidence ($z_3$) & 0.206 & 0.200 \\
    Coordination ($z_5$) & 0.207 & 0.201 \\
    Tool ($z_4$) & 0.218 & 0.212 \\
    Structural ($z_1$) & 0.228 & 0.222 \\
    \bottomrule
  \end{tabular}
\end{table}

Reproduce with \texttt{python -m research.cvec\_dim\_ablation --split test};
outputs under \texttt{results/experiments/cvec5\_dim\_ablation/}.

% Prompt panels + GAIA figure: single-column / full-width (not right-column half width).
% Force single-column appendix pages for prompt panels + GAIA figure (ACL/EMNLP twocolumn).
\clearpage
\onecolumn

% Appendix: QCE planner (decomposition) prompts — verbatim from codebase
\section{QCE Decomposition Prompts}
\label{app:decompose-prompts}

The Query Complexity Estimator (QCE) calls an LLM \emph{procedure planner} before routing.
The planner outputs a JSON object describing an acyclic subtask DAG (Algorithm~\ref{alg:qce}).
Figure~\ref{fig:decompose-prompt-suite} reproduces the released prompts from
\texttt{input/prompts/qce\_decompose.py} without abridgment (default model: \texttt{gpt-4o-mini}).
Agent \emph{execution} prompts are in Appendix~\ref{app:agent-prompts}.

% QCE decomposition prompts — verbatim from input/prompts/qce_decompose.py (full-width figure*)
\begin{figure*}[p]
  \centering
  \begin{tcolorbox}[
    width=0.98\textwidth,
    enhanced,
    colback=white,
    colframe=black!55,
    arc=2pt,
    boxrule=0.6pt,
    title={\sffamily\bfseries QCE Procedure Planner Prompt Suite},
    fonttitle=\small,
    coltitle=white,
    colbacktitle=brown!65!black,
    attach boxed title to top left={yshift=-2mm},
    boxed title style={arc=2pt},
    left=3pt, right=3pt, top=4pt, bottom=3pt,
  ]
  \lstinputlisting[style=promptsuite]{figures/generated/decompose_prompt_suite.txt}
  \end{tcolorbox}
  \caption{QCE decomposition planner prompts reproduced verbatim from
    \texttt{input/prompts/qce\_decompose.py} (system core, per-dataset strategy hints,
    few-shot JSON examples, user template).
    Response is parsed into subtasks, repaired into $G_q$, and passed to QCE feature extraction.}
  \label{fig:decompose-prompt-suite}
\end{figure*}

\clearpage

% Appendix: agent execution prompts — verbatim from codebase
\section{Agent Strategy Prompts}
\label{app:agent-prompts}

At inference the router selects one of four \emph{execution} strategies
($\texttt{raw}$, $\texttt{cot}$, $\texttt{react}$, $\texttt{multiagent}$).
Each strategy uses a dataset-conditioned system prompt from
\texttt{input/prompts/prompts\_core.py}.
Unlike the QCE \emph{procedure planner} (Appendix~\ref{app:decompose-prompts}), these prompts
instruct the model to answer the question under a shared JSON contract
(\texttt{answer}, \texttt{confidence}, \texttt{complexity}) for fair grading across strategies.
Figure~\ref{fig:agent-prompt-suite} reproduces the HotpotQA templates verbatim.

% Agent prompts — HotpotQA verbatim (input/prompts/prompts_core.py); Fig.~6 panel style
\begin{figure*}[p]
  \centering
  \begin{tcolorbox}[
    width=0.98\textwidth,
    enhanced,
    colback=white,
    colframe=black!55,
    arc=2pt,
    boxrule=0.6pt,
    title={\sffamily\bfseries Agent Strategy Prompt Suite (HotpotQA exemplar)},
    fonttitle=\small,
    coltitle=white,
    colbacktitle=brown!65!black,
    attach boxed title to top left={yshift=-2mm},
    boxed title style={arc=2pt},
    left=5pt, right=5pt, top=6pt, bottom=5pt,
  ]
  \scriptsize\sffamily

  \textbf{raw}\par\vspace{3pt}
  {\ttfamily\RaggedRight
  You are a reasoning agent.\\
  Task: Answer the multi-hop question.\\
  Representation rules:\\
  - use ONLY the provided Context paragraphs as evidence\\
  - copy the shortest supporting span from evidence\\
  - do not paraphrase or add words not in evidence\\
  - preserve punctuation, titles, and suffixes\\
  - when multiple candidates appear, pick the shortest span that answers the question\\
  - do not return intermediate hops when the question asks for the final entity\\
  Answer concisely and directly.\\
  Do not explain your reasoning.\\
  Preserve the representation requested by the problem.\\
  Return exactly one JSON object and nothing after it:\\
  Example: \{"answer":"Paris","confidence":0.9,"complexity":0.3\}\\
  answer: short answer span: name, number, date, yes, or no\\
  No markdown fences.\par}

  \vspace{6pt}
  \textbf{cot}\par\vspace{3pt}
  {\ttfamily\RaggedRight
  You are a reasoning agent.\\
  Task: Answer the multi-hop question.\\
  Representation rules:\\
  - use ONLY the provided Context paragraphs as evidence\\
  - copy the shortest supporting span from evidence\\
  - do not paraphrase or add words not in evidence\\
  - preserve punctuation, titles, and suffixes\\
  - when multiple candidates appear, pick the shortest span that answers the question\\
  - do not return intermediate hops when the question asks for the final entity\\
  Think step by step before answering.\\
  Verify intermediate reasoning before the final answer.\\
  Preserve the representation requested by the problem.\\
  Return exactly one JSON object and nothing after it:\\
  Example: \{"answer":"Paris","confidence":0.9,"complexity":0.3\}\\
  answer: short answer span: name, number, date, yes, or no\\
  No markdown fences.\par}

  \vspace{6pt}
  \textbf{react}\par\vspace{3pt}
  {\ttfamily\RaggedRight
  You are a ReAct agent.\\
  Task: Answer the multi-hop question.\\
  Representation rules:\\
  - use ONLY the provided Context paragraphs as evidence\\
  - copy the shortest supporting span from evidence\\
  - do not paraphrase or add words not in evidence\\
  - preserve punctuation, titles, and suffixes\\
  - when multiple candidates appear, pick the shortest span that answers the question\\
  - do not return intermediate hops when the question asks for the final entity\\
  Each completion must be ONE of:\\
  1. Tool step\\
  Thought: (your reasoning)\\
  Action: tool[input]\\
  2. Final step\\
  Thought: (brief justification from Observations)\\
  Final Answer:\\
  (one JSON object --- use a real answer, never template placeholders)\\
  Never generate Observation lines. One Action or one Final Answer per completion --- not both.\\
  Evidence: use ONLY the Context paragraphs below (not the open web).\\
  Use retrieve with short semantic queries over those paragraphs.\\
  Do not answer from outside knowledge.\\
  Example valid tool turn:\\
  Thought: I need evidence about the setting years of Dreamland.\\
  Action: retrieve[Dreamland novel setting years]\\
  Invalid retrieve query (paragraph numbers):\\
  Action: retrieve[1]\\
  Do not use paragraph numbers as retrieve queries.\\
  Preferred format is Thought: then Action:; bare retrieve[query] is accepted only as one-line shorthand.\\
  The runtime supplies Observations after each Action.\\
  Tools:\\
  <tools\_block>\\
  Final Answer:\\
  Return exactly one JSON object and nothing after it:\\
  Example: \{"answer":"Paris","confidence":0.9,"complexity":0.3\}\\
  answer: short answer span: name, number, date, yes, or no\\
  No markdown fences.\\
  Rules:\\
  - use only facts from runtime Observations or Context\\
  - preserve the representation requested by the problem\\
  - Final Answer must be one JSON object with an answer field (see example schema)\\
  - avoid repeating the same or nearly identical tool queries\\
  - the runtime adds Observation lines; never write them yourself\par}

  \vspace{6pt}
  \textbf{multiagent}\par\vspace{3pt}
  {\ttfamily\RaggedRight
  Multi-agent QA: Planner $\rightarrow$ Workers $\rightarrow$ Judge. Strict JSON at each stage.\\
  Task: Answer the multi-hop question.\\
  Representation rules:\\
  - use ONLY the provided Context paragraphs as evidence\\
  - copy the shortest supporting span from evidence\\
  - do not paraphrase or add words not in evidence\\
  - preserve punctuation, titles, and suffixes\\
  - when multiple candidates appear, pick the shortest span that answers the question\\
  - do not return intermediate hops when the question asks for the final entity\\
  Final answer type: short answer span: name, number, date, yes, or no.\par}

  \vspace{6pt}
  \textbf{debate}\par\vspace{3pt}
  {\ttfamily\RaggedRight
  You are a reasoning agent.\\
  Task: Answer the multi-hop question.\\
  Representation rules:\\
  - use ONLY the provided Context paragraphs as evidence\\
  - copy the shortest supporting span from evidence\\
  - do not paraphrase or add words not in evidence\\
  - preserve punctuation, titles, and suffixes\\
  - when multiple candidates appear, pick the shortest span that answers the question\\
  - do not return intermediate hops when the question asks for the final entity\\
  Independent expert --- reason on your own; you do not see other agents.\\
  Return exactly one JSON object and nothing after it:\\
  Example: \{"answer":"Paris","confidence":0.9,"complexity":0.3\}\\
  answer: short answer span: name, number, date, yes, or no\\
  No markdown fences.\par}

  \end{tcolorbox}
  \caption{Agent strategy system prompts for HotpotQA (verbatim from
    \texttt{input/prompts/prompts\_core.py}).
    Routed set $\mathcal{A}=\{\texttt{raw},\texttt{cot},\texttt{react},\texttt{multiagent}\}$;
    \texttt{debate} is the per-expert template inside multi-agent runs.
    Other workloads substitute dataset-specific representation rules and ReAct tool lists.}
  \label{fig:agent-prompt-suite}
\end{figure*}

\clearpage

\clearpage
% Worked GAIA example: planner JSON -> procedure DAG -> C(Q)
\subsection{Worked Example: GAIA Query to Procedure DAG and Complexity Vector}
\label{app:gaia-worked-example}

We walk through a real held-out GAIA transfer query (eval sample
\texttt{f46b4380-207e-4434-820b-f32ce04ae2a4}, Level~2) from decomposition through graph repair to
the five-dimensional complexity vector $\mathbf{z}(q)$ consumed by the router.

\paragraph{Query.}
\begin{quote}\small
It is 1999. Before you party like it is 1999, please assist me in settling a bet.
Fiona Apple and Paula Cole released albums prior to 1999. Of these albums, which didn't receive
a letter grade from Robert Christgau? Provide your answer as a comma delimited list of album titles,
sorted alphabetically.
\end{quote}

\paragraph{Planner output (JSON).}
The procedure planner returns \texttt{final\_constraint=phrase} and four subtasks with
multiple \texttt{depends\_on} edges (no graph repairs on this instance):

\begin{table}[h]
  \centering
  \footnotesize
  \caption{Planner subtasks for the GAIA album-grade query (abridged).}
  \label{tab:gaia-subtasks}
  \begin{tabular}{@{}lllll@{}}
    \toprule
    id & step\_type & tool & depends\_on & goal (abridged) \\
    \midrule
    t1 & retrieve & web\_search & [] & Fiona Apple albums before 1999 \\
    t2 & retrieve & web\_search & [] & Paula Cole albums before 1999 \\
    t3 & retrieve & web\_search & [t1, t2] & Christgau letter grades for albums \\
    t4 & compute & none & [t1, t2, t3] & Albums without letter grade (sink) \\
    \bottomrule
  \end{tabular}
\end{table}

\paragraph{Graph statistics and $\mathbf{z}(q)$.}
QCE computes topology and plan-semantics statistics on $G_q$ (e.g., \texttt{max\_depth}$=3$,
\texttt{n\_nodes}$=4$, multiple merge edges, \texttt{n\_merge\_nodes}${=}2$),
normalizes each statistic with train-fitted caps (\texttt{models/qce\_graph/train\_norm.json}),
and aggregates normalized terms into five groups
(Table~\ref{tab:complexity-dims}; implementation in \texttt{qce/graph.py}).

\begin{table}[h]
  \centering
  \footnotesize
  \caption{Complexity vector for the GAIA album-grade example (train norm applied).}
  \label{tab:gaia-cq}
  \begin{tabular}{@{}lrl@{}}
    \toprule
    Dimension & Value & Main signals (this query) \\
    \midrule
    $z_1$ structural & 0.883 & 4-node DAG, parallel roots, multi-parent sink \\
    $z_2$ reasoning & 0.508 & compute sink after merged retrieval \\
    $z_3$ evidence & 0.719 & three web retrieval hops \\
    $z_4$ tool & 0.486 & repeated \texttt{web\_search}, no attachment \\
    $z_5$ coordination & 0.143 & clean plan; non-trivial merge structure \\
    \bottomrule
  \end{tabular}
\end{table}

\paragraph{Procedure DAG.}
QCE links virtual \texttt{\_\_start\_\_} to both root retrieves, connects each
\texttt{depends\_on} prerequisite, and attaches \texttt{\_\_end\_\_} after the sink
(Figure~\ref{fig:gaia-dag-example}).
The graph is not a simple chain: \texttt{t3} merges \texttt{t1} and \texttt{t2}, and
\texttt{t4} depends on all three predecessors.

% GAIA worked example — top-to-bottom procedure DAG (eval f46b4380)
\begin{figure*}[t]
  \centering
  \resizebox{0.98\textwidth}{!}{%
  \begin{tikzpicture}[
    font=\small\sffamily,
    >=Stealth,
    panel/.style={
      draw=gray!45, rounded corners=4pt, fill=gray!4,
      inner sep=6pt, align=left, font=\scriptsize\sffamily,
    },
    st/.style={
      draw=#1!65!black, fill=#1!12, rounded corners=6pt,
      minimum width=2.0cm, minimum height=1.15cm, align=center,
      line width=0.95pt, inner sep=5pt,
    },
    virtnode/.style={
      draw=gray!50, fill=white, line width=0.75pt,
      rounded corners=4pt, font=\scriptsize\sffamily\ttfamily,
      inner xsep=6pt, inner ysep=3pt,
    },
    dep/.style={->, line width=1.1pt, draw=blue!65!black},
    depside/.style={->, line width=0.85pt, draw=blue!35},
    virtedge/.style={->, line width=0.8pt, draw=gray!55, dashed},
  ]
    % ===== Top metadata =====
    \node[panel, text width=3.85cm] (ctx) at (-4.85, 9.2) {%
      \textbf{Context}\\[2pt]
      1999 bet: Fiona Apple \& Paula Cole albums---which lacked a Christgau \emph{letter} grade? Comma-separated titles (A--Z).
    };
    \node[panel, text width=2.4cm, align=center] (qinfo) at (0, 9.2) {%
      \textbf{Query} $\mathbf{Q}$\\
      GAIA L2 \texttt{f46b4380\ldots}
    };
    \node[panel, text width=3.9cm] (edgebox) at (4.85, 9.2) {%
      \textbf{\texttt{depends\_on}}\\
      \texttt{t1[]}, \texttt{t2[]}; \texttt{t3[t1,t2]}; \texttt{t4[t1,t2,t3]}
    };

    % ===== DAG background (white, behind nodes) =====
    \fill[white, rounded corners=10pt] (-4.05, 0.65) rectangle (4.05, 8.35);
    \draw[gray!45, rounded corners=10pt, line width=0.75pt]
      (-4.05, 0.65) rectangle (4.05, 8.35);

    % ===== DAG nodes: start (top) -> end (bottom) =====
    \node[virtnode] (start) at (0, 8.15) {start};

    \node[st=blue] (t1) at (-2.75, 6.65) {%
      \textbf{t1}\\[2pt]
      {\scriptsize retrieve}\\
      {\scriptsize\texttt{web\_search}}};
    \node[st=blue] (t2) at (2.75, 6.65) {%
      \textbf{t2}\\[2pt]
      {\scriptsize retrieve}\\
      {\scriptsize\texttt{web\_search}}};

    \node[st=blue] (t3) at (0, 4.75) {%
      \textbf{t3}\\[2pt]
      {\scriptsize retrieve}\\
      {\scriptsize\texttt{web\_search}}};

    \node[st=violet] (t4) at (0, 2.65) {%
      \textbf{t4}\\[2pt]
      {\scriptsize compute}\\
      {\scriptsize sink}};

    \node[virtnode] (end) at (0, 1.15) {end};

    % ===== DAG edges (downward flow) =====
    \draw[virtedge] (start.south west) -- (t1.north);
    \draw[virtedge] (start.south east) -- (t2.north);
    \draw[dep] (t1.south) -- (t3.north west);
    \draw[dep] (t2.south) -- (t3.north east);
    \draw[dep] (t3.south) -- (t4.north);
    \draw[virtedge] (t4.south) -- (end.north);
    \draw[depside] (t1.south east) to[out=-58, in=168, looseness=0.85] (t4.north west);
    \draw[depside] (t2.south west) to[out=-122, in=12, looseness=0.85] (t4.north east);

    % ===== Bottom metadata =====
    \node[panel, text width=3.5cm] (leg) at (-3.15, -0.35) {%
      \textbf{Legend}\\
      \tikz{\draw[dep] (0,0) -- (0.55,0);} spine \quad
      \tikz{\draw[depside] (0,0) -- (0.55,0);} direct\\
      \tikz{\draw[virtedge] (0,0) -- (0.55,0);} virtual
    };
    \node[panel, text width=5.2cm, align=left] (zq) at (2.75, -0.35) {%
      \textbf{$\mathbf{z}(q)$} (train norm)\\[3pt]
      \begin{tabular}{@{}l@{\hspace{0.3em}}c@{\hspace{0.2em}}r@{}}
        Structural    & \textcolor{blue!55!black}{\rule{1.5cm}{0.12cm}} & 0.883 \\
        Reasoning     & \textcolor{blue!55!black}{\rule{0.92cm}{0.12cm}} & 0.508 \\
        Evidence      & \textcolor{blue!55!black}{\rule{1.22cm}{0.12cm}} & 0.719 \\
        Tool          & \textcolor{blue!55!black}{\rule{0.88cm}{0.12cm}} & 0.486 \\
        Coordination  & \textcolor{blue!55!black}{\rule{0.20cm}{0.12cm}} & 0.143 \\
      \end{tabular}
    };

  \end{tikzpicture}%
  }

  \caption{Worked GAIA example (eval \texttt{f46b4380-207e-4434-820b-f32ce04ae2a4}): procedure
    DAG $G_q$ (top \texttt{start} $\rightarrow$ bottom \texttt{end}) with parallel retrieves
    \texttt{t1}, \texttt{t2}, merge at \texttt{t3}, and fan-in sink \texttt{t4}
    (Table~\ref{tab:gaia-subtasks}). Curved edges are direct \texttt{depends\_on} links to the sink.
    Normalized $\mathbf{z}(q)$ in Table~\ref{tab:gaia-cq}.}
  \label{fig:gaia-dag-example}
\end{figure*}


The router then concatenates $\mathbf{z}(q)$ with $\psi(q)$ (PCA-16 embedding of the query
text) to form $\phi(q)$ for cost-soft HGBM routing.
This example illustrates how multi-hop GAIA queries with parallel evidence gathering and
merge nodes yield higher structural and evidence complexity than single-file linear plans.


\section{Hyperparameters and Reproducibility}
\label{app:hyperparameters}

\begin{table}[h]
  \centering
  \footnotesize
  \caption{Implementation details and hyperparameters (primary system).}
  \label{tab:hyperparams}
  \begin{tabular}{@{}ll@{}}
    \toprule
    Component & Setting \\
    \midrule
    \multicolumn{2}{@{}l}{\textit{Routing benchmark splits}} \\
    Train / validation / test queries & 808 / 100 / 100 \\
    \midrule
    \multicolumn{2}{@{}l}{\textit{QCE decomposition (Appendix~\ref{app:decompose-prompts})}} \\
    Procedure planner model & \texttt{gpt-4o-mini} \\
    Planner output & JSON subtask DAG $\rightarrow$ $G_q$ \\
    \midrule
    \multicolumn{2}{@{}l}{\textit{Representation (complexity--semantic)}} \\
    Planner complexity dimensions & 5 \\
    Semantic embedding (PCA) & 16 \\
    Total feature dimension & 21 \\
    \midrule
    \multicolumn{2}{@{}l}{\textit{Classifier and supervision}} \\
    Classifier & Gradient-boosted trees (HGBM) \\
    Training target & Cost-soft $p^\star$ (KL loss) \\
    Cost-soft temperature $\tau$ & 100 \\
    Utility-softmax $\tau$ (ablation; val-selected) & 0.05 \\
    \midrule
    \multicolumn{2}{@{}l}{\textit{Evaluation}} \\
    Utility weight $\lambda$ & 25 \\
    Selective gating threshold (transfer) & 0.5 \\
  \end{tabular}
\end{table}

\paragraph{Reproducibility.}
All benchmark splits, oracle outcomes, and feature representations are versioned and
released with the accompanying codebase.
Supplementary figures and per-dataset routing statistics are included in the supplementary
material.

\section{Additional Analyses}
\label{app:additional-analyses}

\paragraph{Routing policy analysis.}
Per-dataset strategy-selection distributions on executed benchmark workloads are provided
in the supplementary material.

\paragraph{Calibration.}
Figure~\ref{fig:calibration} shows reliability diagrams on the held-out test split of the
routing benchmark (generate with
\texttt{python -m research.analysis --split test} $\rightarrow$
\texttt{results/reports/router/test/reliability\_calibration.png}).

\begin{figure}[h]
  \centering
  \IfFileExists{results/reports/router/test/reliability_calibration.png}{%
    \includegraphics[width=0.75\linewidth]{results/reports/router/test/reliability_calibration.png}%
  }{%
    \IfFileExists{results/reports/router/val/reliability_calibration.png}{%
      \includegraphics[width=0.75\linewidth]{results/reports/router/val/reliability_calibration.png}%
    }{%
      \IfFileExists{results/reports/analysis/reliability_calibration.png}{%
        \includegraphics[width=0.75\linewidth]{results/reports/analysis/reliability_calibration.png}%
      }{%
        \fbox{\parbox[c][1in][c]{0.85\linewidth}{\centering\small
          Run \texttt{python -m research.analysis --split test}; expected at
          \texttt{results/reports/router/test/reliability\_calibration.png}.}}%
      }%
    }%
  }
  \caption{Router probability calibration (routing benchmark test split).}
  \label{fig:calibration}
\end{figure}

\paragraph{Cost--accuracy tradeoff.}
Figure~\ref{fig:cost-accuracy} compares cost--accuracy frontiers for routed and fixed
policies on the routing benchmark test split (same analysis run as above, or
\texttt{python -m research.figures} for
\texttt{results/reports/figures/fig3\_cost\_accuracy\_frontier.png}).

\begin{figure}[h]
  \centering
  \IfFileExists{results/reports/router/test/cost_vs_accuracy.png}{%
    \includegraphics[width=0.75\linewidth]{results/reports/router/test/cost_vs_accuracy.png}%
  }{%
    \IfFileExists{results/reports/analysis/cost_vs_accuracy.png}{%
      \includegraphics[width=0.75\linewidth]{results/reports/analysis/cost_vs_accuracy.png}%
    }{%
      \IfFileExists{results/reports/figures/fig3_cost_accuracy_frontier.png}{%
        \includegraphics[width=0.75\linewidth]{results/reports/figures/fig3_cost_accuracy_frontier.png}%
      }{%
        \fbox{\parbox[c][1in][c]{0.85\linewidth}{\centering\small
          Run \texttt{python -m research.analysis --split test}; expected at
          \texttt{results/reports/router/test/cost\_vs\_accuracy.png}.}}%
      }%
    }%
  }
  \caption{Cost--accuracy tradeoff comparing routed and fixed-strategy policies.}
  \label{fig:cost-accuracy}
\end{figure}

\paragraph{Embedding dimension sensitivity.}
We ablate PCA projection size ($8$, $16$, $32$, $64$ dimensions).
Dimension $16$ is used in all primary experiments; larger projections do not improve
held-out regret on the routing benchmark.

\section{Fixed-Strategy Baselines (Per-Agent EM)}
\label{app:fixed-agent-em}

\begin{table}[h]
  \centering
  \footnotesize
  \caption{Fixed-strategy exact match (\%) by workload.}
  \begin{tabular}{@{}lrrrrr@{}}
    \toprule
    Workload & raw & cot & react & multiagent \\
    \midrule
    GAIA & 4.2 & 6.1 & 26.1 & 13.3 \\
    HotpotQA & 39.5 & 36.5 & 69.0 & 32.0 \\
    MuSiQue & 2.5 & 2.5 & 36.5 & 25.5 \\
    MMLU & 48.5 & 53.0 & 59.0 & 46.5 \\
    MATH & 31.0 & 69.5 & 63.0 & 56.0 \\
    \bottomrule
  \end{tabular}
\end{table}

\section{Selective QCE}
\label{app:selective}

An optional selective gate routes queries with high embedding-only confidence without
invoking planner decomposition.
Validation-split regret and executed MMLU-Pro results for this variant are reported in the
supplementary material.

\section{Related Work (Extended)}
\label{app:related-extended}

Section~\ref{sec:related-work} discusses RouteLLM, LLM-Blender, and GraphRouter as the
primary literature baselines for this work.
Table~\ref{tab:related-mapping} summarizes how each system maps to components of our
experimental design.

\begin{table}[h]
  \centering
  \footnotesize
  \caption{Mapping from literature routing baselines to AAR experiments.}
  \label{tab:related-mapping}
  \begin{tabular}{@{}p{0.22\linewidth}p{0.34\linewidth}p{0.36\linewidth}@{}}
    \toprule
    Baseline & Core idea & AAR instantiation \\
    \midrule
    RouteLLM &
      Cost-aware strong/weak model routing &
      Utility $U_a=r_a-\lambda c_a$ and regret (§\ref{sec:experiments-setup}) \\
    LLM-Blender &
      Rank/fuse multiple viable model outputs &
      Cost-soft $p^\star$ over correct strategies (Table~\ref{tab:soft-target}) \\
    GraphRouter &
      Graph-structured query--candidate routing &
      Procedure-DAG complexity $\mathbf{z}(q)$ + embeddings (Table~\ref{tab:feature-regret}) \\
    \bottomrule
  \end{tabular}
\end{table}

RouteLLM~\citep{ong2025routellm}, RouterDC~\citep{chen2024routerdc}, and HybridLLM~\citep{ding2024hybridllm} learn query-conditioned \emph{model} selection under
cost constraints; LLM-Blender~\citep{jiang2023llmblender} operates post-hoc on multiple generations; GraphRouter~\citep{feng2025graphrouter} and
AgentRouter~\citep{zhang2025agentrouter} use entity- or affiliation-centric graphs with GNN routing.
We route among four heterogeneous \emph{agent pipelines} with tabular policies, procedure
graphs, and cost-soft labels restricted to correct strategies.
Direct head-to-head replication of published systems is limited by task mismatch (model/API
routing vs.\ tool-using agent strategies); our tables compare learned and heuristic routers
in the unified four-strategy setting above.

\section{Discussion and Analysis}
\label{app:discussion}

\paragraph{Main result.}
On the routing benchmark, among the evaluated factors cost-soft supervision yields the largest
observed regret reduction: agent regret falls from $0.361$ to $0.192$ relative to hard labels,
and executed exact match rises from 63\% to 82\% (Table~\ref{tab:hard-vs-soft}).
Under cost-soft supervision, the complexity--semantic representation yields the lowest total
regret ($0.198$) and outperforms both embedding-only and heuristic feature sets
(Table~\ref{tab:feature-regret}).
The best configuration combines both components; preserving multiple cost-weighted correct
strategies during training and enriching $\phi(q)$ with planner-derived structure together
improve cost-aware routing under the objective the router is trained to optimize.

\paragraph{Secondary result.}
Benchmark transfer is a secondary evaluation with mixed outcomes: the metric shifts to
executed exact match on heterogeneous public workloads without per-query oracle access
at inference.
The proposed router remains competitive (48.1\% pooled EM) but does not consistently
outperform a simpler heuristic router (50.3\%) or the strongest fixed strategy per family
(Table~\ref{tab:benchmark-transfer-em}); this should not be read as invalidating
routing-benchmark regret gains.
Held-out families (GAIA, MMLU-Pro) and distribution shift likely contribute; the framework
is most reliable where training and evaluation share oracle multi-strategy supervision.
Transfer nevertheless shows that the router does not collapse to a single static policy and
remains usable under full execution.

\paragraph{Oracle reference.}
\textbf{Best agent} performance in Table~\ref{tab:live-benchmark} should be read as an
oracle headroom reference, not as evidence that learned routing is near-optimal on every
workload.
Large gaps between proposed and best-agent exact match (e.g., on GAIA and HotpotQA) indicate
substantial room for improved routing or richer strategy pools, while the routing-benchmark
regret metric captures progress toward the cost-aware objective used during training.

Supplementary analyses---routing policy distributions across workloads, cost--accuracy
tradeoffs on oracle splits, and classifier calibration---are reported in
Appendix~\ref{app:additional-analyses}.
